\documentclass[10pt,a4paper]{amsart}
\usepackage[margin=19mm]{geometry}
\usepackage{amsmath,amssymb,mathtools}
\usepackage[scr=rsfs,cal=euler]{mathalfa}
\usepackage{xfrac}
\usepackage[unicode,hidelinks,bookmarks=false]{hyperref}
\newcommand{\ie}{i.e.\ }
\newcommand{\cf}{cf.\ }
\newcommand{\resp}{respectively\ }
\newcommand{\Subsection}{\subsection}
\providecommand{\nobreakdash}{\nobreak\hbox{-}}
\setlength{\emergencystretch}{2em}
\multlinegap0pt
\usepackage{amssymb}
\usepackage[normalem]{ulem}
\usepackage{tikz}
\usepackage{mathtools}
\newtheorem{lemma}{Lemma}
\title{Levelable graphs}
\author{Kieran Bhaskara, Michael Y. C. Chong, Takayuki Hibi, Naveena Ragunathan, Adam Van Tuyl}
\date{Source: arXiv:2504.02065v4 (CC BY 4.0)}
\begin{document}
\maketitle

\noindent\emph{Source and licence.} Levelable graphs. Version arXiv:2504.02065v4 (CC BY 4.0), distributed by arXiv under the Creative Commons Attribution 4.0 International licence, http://creativecommons.org/licenses/by/4.0/, as recorded in the arXiv licence metadata for this exact version and shown on the abstract page https://arxiv.org/abs/2504.02065v4. Full licence notice: \texttt{licenses/arxiv-2504.02065-CC-BY-4.0.txt}.\par\smallskip

\noindent\textit{Editorial scope.} Counted unit: the article's lemma lemma.connected (source0.tex lines 516--519). The article's complete original proof of it is delivered with it (source0.tex lines 521--572, one \texttt{proof} environment of 52 lines). No mathematics is written, repaired or re-derived: the statement and the proof are the article's own lines, in the article's own notation, and no retained line is substituted or re-flowed.  No further source-local context is needed: every cross-reference of every retained line resolves inside the two delivered spans.  Nothing was omitted from any retained span, so the package carries no omission record.  No retained line contains a citation, so the wrapper carries no bibliography and no bibliography entry.  No retained line contains a figure environment, an image-inclusion command or drawing code, so no image is delivered and the package declares no asset dependency.  Editorial formatting only, in addition: an AMS \texttt{amsart} preamble that restates the article's own declarations and macros used by the retained lines, namely the environments \texttt{lemma} on separate counters, together with the standard packages those lines need.  The retained environments print in this document as Lemma 1 in order of appearance, whereas the article prints the same items with the numbering of its own full document.  The complete native source of the article as distributed in the arXiv e-print for this exact version is bundled unchanged as \texttt{sources/175864/source0.tex}.
\begin{lemma}\label{lemma.connected}
Let $G$ and $H$ be disjoint graphs.  Then $G \cup H$ is levelable if
and only if $G$ and $H$ are levelable.
\end{lemma}

\begin{proof}
Let $G$ be a graph 
on $V(G) = \{x_1,\ldots,x_n\}$ and
$H$ a graph on $V(H) = \{y_1,\ldots,y_m\}$.  
Because the graphs $G$ and $H$ are 
disjoint, we have
\[{\rm MaxInd}(G \cup H) = \{W \cup Y ~|~ W \in {\rm MaxInd}(G),~ 
Y \in {\rm MaxInd}(H)\}.\]

{\bf (``If'')} 
Suppose that $G$ and $H$ are levelable with respect to 
the weight functions $(c_1,\ldots,c_n)$ and 
$(d_1,\ldots,d_m)$, respectively.
Let $c$ be the independence weight of $G$ and $d$ the 
independence weight of $H$.
We claim that $G \cup H$ is a levelable graph
with respect to the weight function
$(e_1,\ldots,e_n,e_{n+1},\ldots,e_{n+m}) = (c_1,\ldots,c_n,d_1,\ldots,d_m)$.
Indeed, for any $Z \in {\rm MaxInd}(G \cup H)$,
we have $Z = W \cup Y$ with $W \cap Y = \emptyset$ and
$W \in {\rm MaxInd}(G)$ and $Y \in {\rm MaxInd}(H)$. Thus
\[\sum_{z_i \in Z} e_i = \sum_{x_i \in W} c_i + \sum_{y_j \in Y} d_j =
c+d.\]

{\bf (``Only If'')} 
Suppose that 
$G \cup H$ is levelable with respect to
the weight function \[(e_1,\ldots,e_n,e_{n+1},\ldots,e_{n+m})\]
with independence weight $e$.
We claim $G$ and $H$ are levelable with respect
to the weight functions
$c = (e_1,\ldots,e_n)$ and 
$d = (e_{n+1},\ldots,e_{n+m})$, respectively.  Fix a
maximal independent set $Y$ of $H$ and consider any
$W_i \in {\rm MaxInd}(G)$.  
Then $W_i \cup Y \in {\rm MaxInd}(G \cup H)$ and so
\[\sum_{z_j \in W_i \cup Y} e_j = \sum_{z_j \in W_i} e_j + 
\sum_{z_j \in Y} e_j = e.\]
Rearranging this equation, we have 
\begin{equation}\label{sum}
\sum_{z_j \in W_i}e_j = e - \sum_{z_j \in Y} e_j.
\end{equation}
Note that since $W_i$ is a subset of $\{x_1,\ldots,x_n\}$, the 
$e_j$'s that appear in $\sum_{z_j \in W_i}e_j$ only appear among
$\{e_1,\ldots,e_n\}$.  Because \eqref{sum} is true for 
all maximal independent sets of $G$, $(e_1,\ldots,e_n)$ is a
weight function that makes $G$ a levelable  
graph with independence weight $e - \sum_{z_j \in Y} e_j$.  

By swapping the roles of $G$ and $H$, the same proof also shows
that $H$ is levelable.
\end{proof}

\end{document}
